%% SCRAP: papers/ACADEMIC_WORDING_GUIDELINES
%% SOURCE: docs/working/papers/ACADEMIC_WORDING_GUIDELINES.md
%% STATUS: CURRENT
%% FITS: ssrn/ch-wording
%% EDITORIAL: lifted — prose rewritten to press voice

\section{Academic Wording Guidelines}

Careful language discipline is essential when a system uses thermodynamic
quantities as conceptual tools rather than as physical claims. This section
codifies the approved vocabulary for all publications and documentation
derived from the StarForth adaptive runtime research, and provides sentence
patterns and review-response templates for common challenge scenarios.

\subsection{The Core Rule: Similarity, Not Equivalence}

Mathematical similarities between StarForth's empirical results and equations
drawn from physics are described with language that asserts structural
resemblance, not physical identity. The distinction is not cosmetic: a claim
of equivalence invites physicists to evaluate the work as a physics paper
(which it is not), while a claim of structural similarity accurately
characterizes the mathematical relationship and places the analogy in the
domain where it belongs---applied statistics and systems modeling.

\subsection{Approved Phrases}

\paragraph{Describing mathematical relationships.}
The preferred formulation evaluates multiple candidate models before selecting
one:

\begin{quote}
``Multiple candidate models were evaluated---linear, polynomial,
exponential, and power-law---and the functional form $f(t) = f_0
e^{-\lambda t}$ provided the best empirical fit with $R^2 = 0.96$.''
\end{quote}

This construction defends against cherry-picking accusations by establishing
that alternatives were tested. Approved phrases for describing the selected
fit include:

\begin{itemize}
  \item \emph{fits the functional form}
  \item \emph{shares structural similarity with}
  \item \emph{is mathematically isomorphic to}
  \item \emph{exhibits the same mathematical structure as}
  \item \emph{the empirical data conforms to}
  \item \emph{matches the pattern of}
\end{itemize}

\paragraph{Describing observations.}
Empirical statements should lead with the measurement, not the model:

\begin{itemize}
  \item \emph{The data reveals\ldots}
  \item \emph{Measurements indicate\ldots}
  \item \emph{The experimental results demonstrate\ldots}
\end{itemize}

\paragraph{Noting physics parallels.}
When drawing a connection to a physics equation or law, use:

\begin{itemize}
  \item \emph{shares dimensional similarity with}
  \item \emph{the functional form also appears in [physics context]}
  \item \emph{this mathematical structure is also found in}
  \item \emph{bears mathematical resemblance to}
\end{itemize}

\subsection{Forbidden Phrases}

The following constructions imply equivalence or causation and must not
appear outside explicitly labelled interpretation sections.

\paragraph{Direct equivalence claims.}
\begin{itemize}
  \item ``is equivalent to'' --- use \emph{fits the functional form}
  \item ``is the same as'' --- use \emph{shares structural similarity with}
  \item ``proves that [physics concept] applies'' --- use \emph{the data is
        consistent with}
\end{itemize}

\paragraph{Unqualified analogies.}
\begin{itemize}
  \item ``is analogous to'' (without the qualifier ``structurally'')
  \item ``corresponds to'' (without ``structurally'' or similar)
  \item ``behaves as'' (without qualification)
  \item ``mirrors'' (implies direct correspondence)
  \item ``exactly as predicted by'' (implies causal application of theory)
\end{itemize}

\paragraph{Causation implications.}
\begin{itemize}
  \item ``because of [physics principle]''
  \item ``due to [physics concept]''
  \item ``governed by [physics law]''
  \item ``obeys [physics equation]''
\end{itemize}

\subsection{Before-and-After Examples}

\paragraph{Example 1: performance scaling.}

\noindent\textbf{Incorrect:}
\begin{quote}
``Performance scales exactly as predicted by special relativity's time
dilation.''
\end{quote}

\noindent\textbf{Correct:}
\begin{quote}
``Multiple candidate performance models were evaluated. The empirical data
fits the functional form $\tau = \tau_0 / \sqrt{1 - \beta^2}$, which shares
structural similarity with the Lorentz transformation from special relativity.''
\end{quote}

The correct version (1) establishes that alternatives were evaluated,
(2) uses ``fits the functional form'' as an empirical observation, and
(3) uses ``shares structural similarity'' as a mathematical fact, not a
causal claim.

\paragraph{Example 2: window capacity.}

\noindent\textbf{Incorrect:}
\begin{quote}
``The system obeys a cosmological constant law $\Lambda(\text{DoF}) =
4096 / (\text{DoF}+1)$.''
\end{quote}

\noindent\textbf{Correct:}
\begin{quote}
``Empirical measurements reveal that window capacity follows the relationship
$\Lambda(\text{DoF}) = 4096 / (\text{DoF}+1)$ with 0.00\% coefficient of
variation. The symbol $\Lambda$ is borrowed from cosmology for mathematical
convenience.''
\end{quote}

The correct version uses ``empirical measurements reveal'' and ``follows the
relationship'' (descriptive, not normative), and explicitly labels $\Lambda$
as notation borrowed for convenience.

\paragraph{Example 3: geodesic convergence.}

\noindent\textbf{Incorrect:}
\begin{quote}
``Workloads converge along geodesics analogous to general relativity.''
\end{quote}

\noindent\textbf{Correct:}
\begin{quote}
``All workload trajectories converge to the same attractor basin regardless of
starting conditions. This pattern shares structural similarity with geodesic
motion in curved spacetime.''
\end{quote}

The correct version states the empirical behavior first, then draws the
mathematical comparison---without implying the physics explains the
computation.

\subsection{Defensive Language Patterns}

Three reusable templates cover the most common situations.

\paragraph{Pattern 1: model evaluation statement.}
\begin{quote}
``[N] candidate models were evaluated, including [list]. The functional form
[equation] provided the best fit with $R^2 = \text{[value]}$.''
\end{quote}

\paragraph{Pattern 2: empirical observation plus mathematical note.}
\begin{quote}
``[Empirical observation]. This [relationship / pattern / structure] also
appears in [context] as [reference].''
\end{quote}

\paragraph{Pattern 3: metaphor disclosure.}
\begin{quote}
``A thermodynamic metaphor is employed, where execution frequency decreases
over time analogously to heat dissipation. The implementation uses exponential
decay $f(t) = f_0 e^{-\lambda t}$ applied at each heartbeat tick.''
\end{quote}

This pattern explicitly labels the metaphor and immediately pairs it with a
literal description of the implementation.

\subsection{Section-Specific Guidelines}

\begin{description}
  \item[Abstracts and titles] Use mathematical or empirical language only.
    \emph{``Frequency-Based Adaptive Runtime''} is acceptable;
    \emph{``Physics-Driven Adaptive Runtime''} is not, in a strict academic context.

  \item[Introductions] State empirical observations before drawing any
    mathematical comparison. Never assert that physics explains computation.

  \item[Methods and implementation sections] Use literal descriptions.
    \emph{``Frequency counter incremented on execution''} is preferred over
    \emph{``temperature increases when word heats up.''} The metaphorical
    vocabulary is reserved for conceptual exposition, not implementation
    description.

  \item[Results sections] Use statistical language. Report
    \emph{``CV = 0.00\% ($p < 10^{-30}$)''} rather than
    \emph{``perfect thermodynamic equilibrium achieved.''}

  \item[Discussion and interpretation sections] Exploratory language is
    permitted with appropriate qualification: \emph{``One interpretation is
    that\ldots however, causation is not established.''} The phrase
    \emph{``this proves that computation follows physics laws''} is never
    acceptable.
\end{description}

\subsection{Reviewer Challenge Responses}

\paragraph{``You are claiming physics.''}
The appropriate response cites the section where the thermodynamic framework
is explicitly introduced as a conceptual metaphor, then points to the
mathematical claim: the functional form $f(t) = f_0 e^{-\lambda t}$ fits
performance data with $R^2 = \text{[value]}$. The fact that this functional
form also appears in physics is noted as a mathematical similarity, not
a physical claim.

\paragraph{``You cherry-picked equations.''}
Cite the model evaluation section documenting that linear, polynomial,
exponential, and power-law candidates were all evaluated, and that the
reported form was selected on empirical fit quality ($R^2$, AIC) rather than
theoretical preference.

\paragraph{``This is pseudoscience.''}
All claims are empirically verifiable. The physics references describe
mathematical similarities, documented explicitly as metaphorical framing
in the ontology section. Point the reviewer to the formal claim table
and the reproduction protocol.

\subsection{Pre-Publication Checklist}

Before submitting any paper, presentation, or externally visible documentation,
verify the following substitutions have been applied throughout:

\begin{itemize}
  \item Replace \emph{``is equivalent to''} with \emph{``fits the functional form''}
  \item Replace \emph{``is analogous to''} with \emph{``shares structural similarity with''}
  \item Add the qualifier \emph{``structurally''} to any unqualified \emph{``corresponds to''}
  \item Replace \emph{``mirrors''} with \emph{``matches the pattern of''}
  \item Replace \emph{``exactly as''} with \emph{``according to the functional form''}
  \item Replace \emph{``obeys''} with \emph{``follows the relationship''}
  \item Replace \emph{``governed by''} with \emph{``characterized by''}
  \item Verify that all metaphors are explicitly labelled in running text
  \item Confirm that all empirical claims carry data references
  \item Confirm that all physics parallels use similarity language
\end{itemize}

\subsection{The Golden Rule}

Describe what was \emph{observed}; note where the observations
\emph{mathematically resemble} known equations; do not claim the physics
\emph{causes} the behavior.

The defensible position is: ``Measurements in this system fit these
mathematical forms. These forms also appear in physics. This similarity may
provide useful modeling frameworks.'' The indefensible position is: ``Our
system implements physics'' or ``physics laws govern our system.''
